\documentclass[final,a4paper,10pt,onecolumn]{article}
\usepackage[polish]{cv}

\begin{document}

\cvheader{Maciej}{Flaga}{mflaga@student.agh.edu.pl}{+48 576 396 580}{Rząska, Polska}

% ---------- Sidebar ----------
\begin{cvsidebar}

\section{Oprogramowanie}

\skillgroup{Programowanie i HPC}{
    \imgitem{img/soft/cpp.png}{C/C++}
    \imgitem{img/soft/cuda.png}{CUDA C++}
}
\skillgroup{Analiza danych}{
    \imgitem{img/soft/py.png}{Python}
    \imgitem{img/soft/sql.png}{Podstawy PostgreSQL}
}
\skillgroup{Symulacje / Modelowanie}{
    \imgitem{img/soft/fusion.png}{Autodesk Fusion}
    \imgitem{img/soft/comsol.jpeg}{COMSOL}
}
\skillgroup{Narzędzia deweloperskie}{
    \imgitem{img/soft/linux.png}{Linux}
    \imgitem{img/soft/git.png}{Git}
    \imgitem{img/soft/anthropic.png}{Claude Code}
}
\skillgroup{Skład tekstu}{
    \imgitem{img/soft/latex2.png}{\LaTeX}
}

\section{Języki}

\begin{iconlist}
    \flagitem{img/flag/pl.png}{Polski -- ojczysty}
    \flagitem{img/flag/gb.png}{Angielski -- biegły}
    \flagitem{img/flag/de.png}{Niemiecki -- podstawowy}
\end{iconlist}

\section{Certyfikaty}

\begin{iconlist}
    \imgitem{img/soft/acert.png}{ACERT Business English C1}
\end{iconlist}

\section{Zainteresowania}

\interestgroup{Naukowe}{
    \item Cyfrowe modelowanie rzeczywistości
}
\interestgroup{Pozanaukowe}{
    \item Wspinaczka
    \item Turystyka górska
    \item Druk 3D i grafika 3D
}

\end{cvsidebar}
% ---------- Main column ----------
\begin{cvmain}

\section{Profil}

Student Nanoinżynierii Materiałów o szczególnym zainteresowaniu fizyką obliczeniową, symulacjami numerycznymi oraz teoretycznymi aspektami inżynierii materiałowej.
Doświadczony w wykorzystaniu C/C++, Pythona i CUDA w obliczeniach naukowych. 
Chcę rozwijać umiejętności w zakresie symulacji, modelowania i analizy danych.

\section{Wykształcenie}

\eduentry
    {img/agh.png}
    {Student}
    {Akademia Górniczo-Hutnicza w Krakowie}
    {Kraków, Polska}
    {paź. 2023 -- obecnie}
    {Wydział Fizyki i Informatyki Stosowanej\\
    Wydział Inżynierii Materiałowej i Ceramiki\\
    Kierunek: Nanoinżynieria Materiałów\\
    Specjalność: Modelowanie}

\section{Doświadczenie zawodowe}

\jobentry
    {Analityk danych}
    {Instytut Fizyki Jądrowej im. Henryka Niewodniczańskiego PAN}
    {Kraków, Polska}
    {lip. 2026 -- wrz. 2026 (3 mies.)}
    {Analiza wieloletnich danych z monitorów neutronów i katalogów trzęsień ziemi w Pythonie, z wykorzystaniem metod Monte Carlo do badania korelacji między promieniowaniem kosmicznym a aktywnością sejsmiczną.}

\jobentry
    {Pracownik hali produkcyjnej}
    {YK Poland Sp. z o.o.}
    {Modlniczka, Polska}
    {lip. 2024 -- wrz. 2024 (3 mies.)}
    {Realizacja pełnego cyklu produkcji ramek metalowych: od cięcia profili i montażu elementów, po foliowanie termiczne oraz finalne pakowanie.}

\jobentry
    {Pracownik hali produkcyjnej}
    {BMK Professional Electronics}
    {Augsburg, Niemcy}
    {lip. 2023 -- wrz. 2023 (3 mies.)}
    {Montaż komponentów na płytkach PCB, obsługa procesu lutowania w piecu przelotowym oraz przygotowanie obwodów do produkcji poprzez maskowanie i maszynowe nakładanie powłok ochronnych.}

\section{Projekty akademickie}

\projentry
    {\href{https://github.com/mflagga/cuda-cfd-solver}{Symulacja transportu masy}}
    {
        Stworzenie akcelerowanego na GPU silnika obliczeniowej dynamiki płynów (CFD) do symulacji przepływu nieściśliwego i transportu masy w oparciu o równania \textbf{Naviera-Stokesa} oraz \textbf{adwekcji-dyfuzji}.
        Implementacja sformułowania funkcja prądu -- wirowość metodą różnic skończonych; powstałe układy równań liniowych rozwiązywane są \textbf{równoległą relaksacją Gaussa-Seidela (red-black)} we własnych kernelach \textbf{CUDA}.
        Symulacja wykorzystuje \textbf{schemat Cranka-Nicolson} do stabilnego całkowania rozkładu masy w czasie oraz potok \textbf{Python/FFmpeg} do renderowania animacji pól prędkości i map stężenia.
    }

\projentry
    {\href{https://github.com/mflagga/wave-vis}{Symulacja równania falowego (2+1)D}}
    {
        Stworzenie wydajnego wizualizatora równania falowego (2+1)D opartego na dokładnym rozwiązaniu analitycznym wyprowadzonym metodą \textbf{separacji zmiennych} i \textbf{rozwinięcia w szereg Fouriera}.
        Numeryczne obliczanie podwójnego szeregu nieskończonego w \textbf{CUDA C++} z równoległym sumowaniem na GPU, co umożliwia obliczenia na siatkach o wysokiej rozdzielczości.
        Projekt obejmuje w pełni zautomatyzowany potok łączący backend w C++, \textbf{Python (Matplotlib)} do wykresów powierzchni 3D oraz \textbf{FFmpeg} do renderowania wideo, sterowany własnym \textbf{Makefile}.
    }

\end{cvmain}

\end{document}
